\chapter{ Conclusion }

This project documents investigation of WebAssembly (WASM) as a hardware-agnostic intermediate representation for distributing audio DSP modules to embedded targets. By benchmarking three execution strategies on a representative embedded platform and creating proof-of-concept developer tooling, I have evaluated the practicality of WASM-based processing in resource-constrained audio systems and presented an example workflow for development and distribution.

Performance measurements indicate that interpreter-based execution is not suitable for real-time embedded DSP. However, approaches leveraging ahead-of-time (AOT) translation mechanisms demonstrate acceptable performance for a meaningful subset of audio algorithms.

These findings support the broader claim that a carefully specified WASM ABI for audio plugins could form the basis of a universal DSP distribution format, similar to that proposed in \cite{gaster2020audio}. Such a format would allow algorithms written in multiple programming languages to compile to a single intermediate representation and be uniformly deployed across embedded devices, desktop applications, and web platforms using appropriate runtimes.

The primary contributions of this work are: (1) the exploration of multiple execution strategies for WASM-based audio DSP modules, (2) empirical benchmarking of their relative performance on representative embedded hardware, and (3) creation of web-based developer tooling to facilitate their authoring and deployment. Future work will include refinement of a WASM ABI for audio plugins, further exploration of compiler and runtime optimizations, and evaluation on additional deployment targets.
